\documentclass[11pt]{article}
\usepackage{amsmath,amssymb,amsfonts,mathtools}
%\usepackage{hyperref}
\usepackage[margin=1in]{geometry}

\newcommand{\tr}{\mathrm{tr}}
\newcommand{\ii}{\mathrm{i}}
\newcommand{\dd}{\mathrm{d}}
\newcommand{\eps}{\varepsilon}
\usepackage{jheppub} % for details on the use of the package, please see the JINST-\
\newcommand{\nn}{\nonumber}
\newcommand{\lb}{\left(}
\newcommand{\rb}{\right)}
\usepackage{amsmath, amssymb}

\newcommand{\NL}[1]{\textcolor{blue}{#1}}


\begin{document}

\section{Relative entropy for mixed excited states in Free Boson CFT}

\subsection*{Setup and Replica Trick}

We consider a massless free boson conformal field theory (CFT) in $1+1$ dimensions on a cylinder of circumference $R$. We focus on a spatial subsystem $A$ defined by the interval $(-l/2, l/2)$ at a fixed time slice.

Our states of interest are the reduced density matrices corresponding to global excited coherent states. These states are generated by the action of vertex operators on the vacuum at past infinity:
\begin{equation}
|\alpha\rangle = V_{\alpha}|\Omega\rangle = :e^{i\alpha\phi}:|\Omega\rangle\,.
\end{equation}
The reduced density matrix of the state $|\alpha\rangle$ on the subsystem $A$ is denoted by $\rho_\alpha$.

We aim to compute the relative entropy $S(\rho_\alpha \| \sigma_y)$ where the reference state is the mixture
\begin{equation}
\sigma_y = y\rho_{\alpha} +(1-y)\rho_{\beta},\qquad 0\leq y\leq 1\,.
\end{equation}
The standard replica trick for relative entropy relates this quantity to the analytic continuation of trace functionals that break replica symmetry. Specifically, for the mixed state case, we utilize the following identity derived from the generalized replica trick:
\begin{equation}
S\left(\rho_\alpha \Big\| y\rho_\alpha+(1-y)\rho_\beta\right)
=
\left.
\partial_{n} \log \left[
\frac{\tr(\rho_\alpha^{n}) \,[\tr(\sigma_y)]^{n-1}}
{\tr(\rho_\alpha\,\sigma_y^{n-1}) \,[\tr(\rho_\alpha)]^{n-1}}
\right]\right|_{n\to 1}\,.
\label{eq:replica-identity}
\end{equation}
The final answer will be \eqref{eq:S-final}.

The most non-trivial component of this calculation is the term $\tr(\rho_\alpha\,\sigma_y^{n-1})$. Expanding this using the binomial theorem:
\begin{eqnarray}
\tr\!\big(\rho_\alpha (y\rho_\alpha+(1-y)\rho_\beta)^{n-1}\big)
&&= \sum_{p=0}^{n-1} \binom{n-1}{p} y^{n-1-p}(1-y)^p\,
\tr(\rho_{\alpha}^{\,n-p} \rho_{\beta}^p)
+\text{terms with commutators}\,.
\label{eq:binomial-expansion}
\end{eqnarray}
Using the operator-state correspondence, the trace $\tr(\rho_{\alpha}^{n-p} \rho_{\beta}^p)$ corresponds to a path integral on an $n$-sheeted Riemann surface with specific operator insertions. In the limit $n \to 1$, we can identify this with the correlation function of $p$ operators of type $\beta$ and $q=n-p$ operators of type $\alpha$:
\begin{equation}\label{fullcorr}
C_{q,p} \equiv \langle \mathcal{O}_{\alpha}^{(q)} \mathcal{O}_{\beta}^{(p)} \rangle\,.
\end{equation}
As was pointed out, in the case of free bosons, the order of operators will not matter in this correlator, and the commutator terms vanish.

\subsection*{Operator Insertions on the Strip}

We perform the calculation on a strip of width $2\pi$, obtained via the standard conformal mapping from the $n$-sheeted cylinder. The replica sheets are indexed by $k = 0, \dots, n-1$.

The operator insertion points $t_k^{\pm}$ on the strip are:
\begin{equation}
t_k^{\pm} = \frac{2\pi k}{n} \pm \frac{\pi x}{n}\,,
\qquad x = \frac{l}{R}\,,
\end{equation}
where $x$ is the dimensionless interval size. For the correlator $\langle \mathcal{O}_{\alpha}^{(q)} \mathcal{O}_{\beta}^{(p)} \rangle$ with $q=n-p$, the operators are explicitly distributed as follows:
on sheets $k=0, \dots, q-1$ ($\alpha$-sector) we insert $V_{\alpha}$ at $t_k^+$ and $V_{-\alpha}$ at $t_k^-$, and on sheets $k=q, \dots, n-1$ ($\beta$-sector) we insert $V_{\beta}$ at $t_k^+$ and $V_{-\beta}$ at $t_k^-$.

\subsection*{Calculation of the correlator}

The correlation function for primary vertex operators $V_{Q} = :e^{i Q \phi}:$ is given by the product of distances raised to the power of the product of their charges. The chordal distance on the strip is
\begin{equation}
d(t_i, t_j) = 2 \sin\left(\frac{t_i - t_j}{2}\right)\,.
\end{equation}
Hence
\begin{equation}
\left\langle \prod_{i} V_{Q_i}(t_i) \right\rangle \propto
\prod_{i<j} \left[ 2 \sin\left( \frac{t_i - t_j}{2} \right) \right]^{Q_i Q_j}\,.
\end{equation}
\begin{itemize}
    \item \textbf{Like charges} ($Q_i Q_j > 0$): the exponent is $+Q_i Q_j$.
    \item \textbf{Opposite charges} ($Q_i Q_j < 0$): the exponent is $-|Q_i Q_j|$.
\end{itemize}

We have the following contributions.

\paragraph{1. Nearest Neighbor Interactions (Same sheet, same sector).}
These are the contractions between $t_k^+$ and $t_k^-$ on the same sheet.
\begin{itemize}
    \item $\alpha$-sector ($q$ pairs): $(\alpha, -\alpha)$ separated by $\Delta t = 2\pi x/n$.
    \[
    \implies \left( 2 \sin \frac{\pi x}{n} \right)^{-q \alpha^2}.
    \]
    \item $\beta$-sector ($p$ pairs): $(\beta, -\beta)$ separated by $\Delta t = 2\pi x/n$.
    \[
    \implies \left( 2 \sin \frac{\pi x}{n} \right)^{-p \beta^2}.
    \]
\end{itemize}
Putting the two together, we have
\begin{equation}\label{Nfunc}
  \mathcal{N}(q,p;x)=\left( 2 \sin \frac{\pi x}{n} \right)^{-(q\alpha^2+p\beta^2)}\,.
\end{equation}

\paragraph{2. Intra-sector Interactions (Different sheets, same sector).}
For the $\alpha$-sector, we sum over sheet separations $k = 1, \dots, q-1$. The number of pairs at distance $k$ is $(q-k)$.
\begin{itemize}
    \item Like charges $(\alpha, \alpha)$ and $(-\alpha, -\alpha)$:
    \[
    \text{Dist} = \frac{2\pi k}{n}
    \qquad\Rightarrow\qquad
    \text{net exponent } = +2\alpha^2.
    \]
    \item Opposite charges $(\alpha, -\alpha)$: two diagonals $\frac{2\pi k}{n} \pm \frac{2\pi x}{n}$.
    \[
    \text{Dist} = \frac{2\pi (k\pm x)}{n}
    \qquad\Rightarrow\qquad
    \text{exponent } = -\alpha^2 \text{ (each)}.
    \]
\end{itemize}
Collecting these terms for the $\alpha$ sector gives $[g(q;x)]^{\alpha^2}$:
\begin{equation}
g(q;x) = \prod_{k=1}^{q-1} \left( \frac{\sin^2(\frac{\pi k}{n})}{\sin(\frac{\pi(k+x)}{n}) \sin(\frac{\pi(k-x)}{n})} \right)^{q-k}\,.
\end{equation}
The $\beta$ sector is identical with $q \to p$ and $\alpha \to \beta$, and we have the overall contribution
\begin{equation}
g(q;x)^{\alpha^2}\,g(p;x)^{\beta^2}\,.
\end{equation}

\paragraph{3. Cross-sector Interactions.}
Interactions between the $q$ $\alpha$-sheets and $p$ $\beta$-sheets yield terms involving the product $\alpha \beta$. These collect into $[h(q,p;x)]^{\alpha\beta}$:
\begin{align}\label{hfunct}
h(q,p;x)
&= \prod_{l=1}^{p} \prod_{k=1}^{q}
\frac{\sin^2(\frac{\pi(k+l-1)}{n})}{\sin(\frac{\pi(k+l-1+x)}{n}) \sin(\frac{\pi(k+l-1-x)}{n})}\,.
\end{align}

Putting all the pieces together, we find that the full correlator in \eqref{fullcorr} is
\begin{equation}
\label{eq:mixed_correlator}
\langle \mathcal{O}_{\alpha}^{(q)} \mathcal{O}_{\beta}^{(p)} \rangle =
\mathcal{N}(q,p;x)\,
\left[ g(q;x) \right]^{\alpha^2}
\left[ g(p;x) \right]^{\beta^2}
\left[ h(q,p;x) \right]^{\alpha\beta}\,.
\end{equation}

We need to analytically continue in $n$. Define the replica sum
\begin{eqnarray}\label{Fn}
\tr(\rho_\alpha\,\sigma_y^{\,n-1})
&&= \sum_{p=0}^{n-1} \binom{n-1}{p}\,y^{n-1-p}(1-y)^p\,
\langle \mathcal{O}_{\alpha}^{(n-p)} \mathcal{O}_{\beta}^{(p)} \rangle\,.
\end{eqnarray}

\subsection*{Analytic continuation in $n$}

We consider the relative entropy $S(\rho_\alpha \| \sigma_y)$ where $\sigma_y = y \rho_\alpha + (1-y)\rho_\beta$ for excited coherent states in a $1+1$D free boson CFT. The replica trick requires evaluating:
\begin{equation}
\tr(\rho_\alpha \sigma_y^{n-1}) = \sum_{p=0}^{n-1} \binom{n-1}{p} y^{n-1-p} (1-y)^p\, C_{n-p, p},
\end{equation}
where $C_{n-p, p} = \langle \mathcal{O}_\alpha^{(n-p)} \mathcal{O}_\beta^{(p)} \rangle$ is the mixed correlator on the $n$-sheeted replica surface.

\paragraph{Step 1: Log-Sum expansion of the correlator.}
The correlator is a product of geometric functions $g$ and $h$. We focus on the intra-sector term $\log g(p;x)$. For integer $p$, it is defined as:
\begin{equation}
\log g(p;x) = \sum_{k=1}^{p-1} (p-k) \left[ 2 \log d\left(\frac{k}{n}\right) - \log d\left(\frac{k+x}{n}\right) - \log d\left(\frac{k-x}{n}\right) \right],
\end{equation}
where $d(y) = 2\sin(\pi y)$ is the distance between operator insertions. We utilize the Fourier expansion for $\log \sin(\pi y)$:
\begin{equation}
\log \sin(\pi y) = -\log 2 - \sum_{m=1}^{\infty} \frac{\cos(2\pi m y)}{m}\,.
\end{equation}
Substituting this into the potential $\mathcal{V}(k,x)$ inside the sum, the constants cancel, yielding:
\begin{equation}
\mathcal{V}(k,x) = -2 \sum_{m=1}^{\infty} \frac{1}{m} \left[ 1 - \cos\left(\frac{2\pi m x}{n}\right) \right] \cos\left(\frac{2\pi m k}{n}\right).
\end{equation}
The expression above has the advantage that it has separated the contribution of $k$ from those of $x$, and we can now perform the sum over $k$.

\paragraph{Step 2: Sum over $k$.}
Following Section 6 of Calabrese, Cardy, and Tonni (2011), we perform the sum over $k$ for each mode $m$:
\begin{equation}
\log g(p;x) = -2 \sum_{m=1}^{\infty} \frac{1 - \cos(2\pi m x/n)}{m}\, \sum_{k=1}^{p-1} (p-k) \cos\left(\frac{2\pi m k}{n}\right).
\end{equation}
The inner sum $S_m(p) = \sum_{k=1}^{p-1} (p-k) \cos(k\theta_m)$ is evaluated using geometric series. For integer $p$ and $\theta_m=2\pi m/n$:
\begin{equation}
S_m(p) =
\mathrm{Re} \left[
\frac{p e^{\ii\theta_m}}{1-e^{\ii\theta_m}} -
\frac{e^{\ii\theta_m}(1-e^{\ii p\theta_m})}{(1-e^{\ii\theta_m})^2}
\right]
=
-\frac{p}{2}+\frac{1-\cos(p\theta_m)}{4\sin^2(\theta_m/2)}\,.
\end{equation}
Therefore, we find
\begin{eqnarray}
\log g(p;x)
&=&
\sum_{m=1}^\infty \frac{1-\cos(\frac{2\pi m x}{n})}{m}
\lb p-\frac{\sin^2(\pi p m/n)}{\sin^2(\pi m/n)}\rb\,.
\end{eqnarray}

Define the mode kernels
\begin{equation}
L(n,x):=\sum_{m=1}^\infty \frac{1-\cos(2\pi m x/n)}{m},
\qquad
K_m(n,x):=\frac{1-\cos(2\pi m x/n)}{2m\,\sin^2(\pi m/n)}.
\label{eq:L-Km}
\end{equation}
to write
\begin{equation}
\log g(p;x)= p\,L(n,x)\;-\;\sum_{m=1}^\infty K_m(n,x)\,\bigl(1-\cos(p\theta_m)\bigr),
\qquad \theta_m=\frac{2\pi m}{n}.
\label{eq:logg-final}
\end{equation}
All $p$-independent pieces can be absorbed into $\Omega_{n}^{\rm (rest)}(x)$ later.

A technical note to keep in mind is that to make every rearrangement fully justified, one may define all mode sums with a common Abel factor
$e^{-\eps m}$ in the summand and take $\eps\to 0^+$ at the end. This guarantees absolute and uniform convergence in intermediate steps.

Next, we move to the function $h(q,p;x)$ in \eqref{hfunct} which includes a double product over $(k,l)$ with $1\le k\le q$, $1\le l\le p$.
Taking logs and using the same log-sum representation as we did for the $g$ function yields
\begin{equation}
\log h(q,p;x)=\sum_{r=1}^{n-1} w_r(p,q)\,V(r,x),
\label{eq:logh-weighted}
\end{equation}
where $w_r(p,q)$ counts the number of pairs $(k,l)$ such that $k+l-1=r$.

It is convenient to evaluate the weighted cosine sum via generating functions.
Let $\theta$ be arbitrary and note
\begin{equation}
\sum_{r=1}^{n-1} w_r(p,q)e^{\ii r\theta}
=\sum_{k=1}^{q}\sum_{l=1}^{p} e^{\ii(k+l-1)\theta}
=e^{-\ii\theta}\Big(\sum_{k=1}^{q}e^{\ii k\theta}\Big)\Big(\sum_{l=1}^{p}e^{\ii l\theta}\Big).
\label{eq:w-gen}
\end{equation}
Now specialize to the case $q=n-p$ and $\theta=\theta_m=2\pi m/n$.
Using $\sum_{k=1}^{n-p} e^{\ii k\theta_m}=-e^{-\ii p\theta_m/2}\,\frac{\sin(p\theta_m/2)}{\sin(\theta_m/2)}$
and $\sum_{l=1}^{p} e^{\ii l\theta_m}=e^{\ii(p+1)\theta_m/2}\,\frac{\sin(p\theta_m/2)}{\sin(\theta_m/2)}$,
the right-hand side of \eqref{eq:w-gen} becomes real and equals
\begin{equation}
\sum_{r=1}^{n-1} w_r(p,n-p)\cos(r\theta_m)
= -\left(\frac{\sin(p\theta_m/2)}{\sin(\theta_m/2)}\right)^2
= -\frac{1-\cos(p\theta_m)}{2\sin^2(\theta_m/2)}.
\label{eq:wcos-eval}
\end{equation}
Substituting \eqref{eq:wcos-eval} into \eqref{eq:logh-weighted} gives the exact mode form
\begin{equation}
\log h(n-p,p;x)= 2\sum_{m=1}^\infty K_m(n,x)\,\bigl(1-\cos(p\theta_m)\bigr),
\label{eq:logh-final}
\end{equation}
with the \emph{same} $K_m(n,x)$ as in \eqref{eq:L-Km}.

We are now ready to put together all the pieces back into $C_{n-p,p}$.
Using \eqref{eq:logg-final} and \eqref{eq:logh-final}, we find
\begin{align}
\alpha^2\log g(q;x)+\beta^2\log g(p;x)+\alpha\beta\log h(q,p;x)
&=(\alpha^2 q+\beta^2 p)\,L(n,x)
\nonumber\\
&\quad-(\alpha-\beta)^2\sum_{m=1}^\infty K_m(n,x)\,\bigl(1-\cos(p\theta_m)\bigr),
\label{eq:combo}
\end{align}
where we have included the nearest-neighbor factor $\mathcal{N}$ into the definition of $\Lambda$ below.
Absorbing all $p$-independent constants (including the $-(\alpha-\beta)^2\sum_m K_m$ term) into
$\Omega_n(x)$, we obtain
\begin{equation}
C_{n-p,p}(x)
=\Omega_n(x)\,\Lambda_\alpha^{\,n-p}\,\Lambda_\beta^{\,p}\,
\exp\!\left[\sum_{m=1}^\infty a_m(n,x)\cos(p\theta_m)\right],
\label{eq:C-expcos}
\end{equation}
where
\begin{equation}
a_m(n,x)=(\alpha-\beta)^2\,K_m(n,x)
=(\alpha-\beta)^2\,\frac{1-\cos(2\pi m x/n)}{2m\,\sin^2(\pi m/n)}\ .
\label{eq:am}
\end{equation}

\subsection*{Fourier/Bessel expansion and definition of the $R$-modes}

Let $\theta:=\frac{2\pi p}{n}$ so that $\cos(p\theta_m)=\cos(m\theta)$.
Using $e^{a\cos(m\theta)}=\sum_{r\in\mathbb{Z}} I_r(a)\,e^{\ii r m\theta}$ (modified Bessel functions),
we expand
\begin{equation}
\exp\!\left[\sum_{m\ge 1} a_m\cos(m\theta)\right]
=\sum_{R\in\mathbb{Z}} W_R(n,x;\alpha,\beta)\,e^{\ii R\theta},
\label{eq:WR-Fourier}
\end{equation}
where the coefficients are the (infinite) Bessel convolution
\begin{equation}
W_R(n,x;\alpha,\beta)=
\sum_{\{r_m\}\,:\,\sum_{m\ge1} m r_m=R}\;\prod_{m\ge1} I_{r_m}\!\big(a_m(n,x)\big),
\qquad
W_{-R}=W_R.
\label{eq:WR-def}
\end{equation}
This allows us to rewrite
\begin{equation}
C_{n-p,p}(x)
=\Omega_n\,\Lambda_\alpha^{\,n-p}\Lambda_\beta^{\,p}\;
\sum_{R\in\mathbb{Z}} W_R(n)\,e^{\ii \frac{2\pi R}{n}p}.
\label{eq:C-R}
\end{equation}

The expression in \eqref{eq:C-R} has the advantage that if we can swap the sum over $R$ with the sum over $p$, we can perform the $p$-sum explicitly. This is our next step. With an Abel regulator, the Fourier coefficients $W_R$ decay exponentially and in particular satisfy $\sum_R |W_R|<\infty$, so swapping the sums is justified by Tonelli/Fubini.

Plugging \eqref{eq:C-R} back into \eqref{Fn} and swapping the order of summation, we obtain
\begin{align}
\tr(\rho_\alpha\,\sigma_y^{\,n-1})
&=\Omega_n\,\Lambda_\alpha\sum_{R\in\mathbb{Z}}W_R(n)
\sum_{p=0}^{n-1}\binom{n-1}{p}\,
(y\Lambda_\alpha)^{n-1-p}\,
\bigl((1-y)\Lambda_\beta e^{\ii2\pi R/n}\bigr)^{p}
\nonumber\\
&=
\Omega_n\,\Lambda_\alpha\sum_{R\in\mathbb{Z}}W_R(n)\,
\Bigl(y\Lambda_\alpha+(1-y)\Lambda_\beta e^{\ii2\pi R/n}\Bigr)^{n-1}\ .
\label{eq:F-R}
\end{align}

Next, we would like to argue that only the $R=0$ term in the expression above contributes to the $n\to 1$ limit at the order relevant for relative entropy.
To establish this, define $\eps=n-1$ and
\begin{equation}
A_R(n):=y\Lambda_\alpha+(1-y)\Lambda_\beta e^{\ii2\pi R/n},\qquad A_0:=y\Lambda_\alpha+(1-y)\Lambda_\beta>0.
\end{equation}
Then \eqref{eq:F-R} reads
\begin{equation}
\tr(\rho_\alpha\,\sigma_y^{\,n-1})=\Omega_n\Lambda_\alpha\sum_{R\in\mathbb{Z}}W_R(n)\,A_R(n)^{\eps}.
\label{eq:F-eps}
\end{equation}

With an Abel regulator, $W_R$ decays exponentially and has finite moments.
Expand the phase for $n=1+\eps$:
\begin{equation}
e^{\ii2\pi R/n}=e^{\ii2\pi R/(1+\eps)}=\exp\!\bigl(-\ii2\pi R\eps+O(R\eps^2)\bigr),
\end{equation}
so $A_R(n)=A_0+\delta_R$ with $\delta_R=O(R\eps)$ for fixed $R$ and
$\delta_{-R}=\overline{\delta_R}$. Using
\begin{equation}
A_R^{\eps}=\exp\!\bigl(\eps\log(A_0+\delta_R)\bigr)=A_0^{\eps}\left[1+\eps\frac{\delta_R}{A_0}+O(\eps|\delta_R|^2)\right],
\end{equation}
and the evenness $W_{-R}=W_R$, the potentially dangerous linear term
$\sum_R W_R\,\delta_R$ cancels at $O(\eps)$ (its leading part is odd/purely imaginary),
leaving only $O(\eps^2)$ corrections after summing over $R$.
Consequently, after absorbing $\sum_R W_R$ into $\Omega_n$ (a $p$-independent normalization),
\begin{equation}
\tr(\rho_\alpha\,\sigma_y^{\,n-1})
=\Omega_{1+\eps}\Lambda_\alpha\,A_0^{\eps}\bigl(1+O(\eps^2)\bigr).
\label{eq:F-asymp}
\end{equation}
Thus the $R\neq 0$ modes do not affect $\partial_n\log \tr(\rho_\alpha\,\sigma_y^{\,n-1})$ at $n=1$.

The replica prescription for the relative entropy is
\begin{equation}
S(\rho_\alpha\|\sigma_y)=
\left.\partial_n\Bigl(\log \tr(\rho_\alpha^n)-\log \tr(\rho_\alpha\sigma_y^{n-1})\Bigr)\right|_{n=1}.
\label{eq:S-replica}
\end{equation}
In the same framework one has $\tr(\rho_\alpha^n)=\widetilde\Omega_n\,\Lambda_\alpha^{\,n}$ so that
$\partial_n\log \tr(\rho_\alpha^n)|_{n=1}=\log\Lambda_\alpha$ (all $\widetilde\Omega_n$ terms cancel
against $\Omega_n$ in the difference).
Using \eqref{eq:F-asymp},
\begin{equation}
\left.\partial_n\log \tr(\rho_\alpha\sigma_y^{n-1})\right|_{n=1}
=\log A_0=\log\!\bigl(y\Lambda_\alpha+(1-y)\Lambda_\beta\bigr).
\end{equation}
Therefore
\begin{equation}
S(\rho_\alpha\|\sigma_y)=\log\Lambda_\alpha-\log\!\bigl(y\Lambda_\alpha+(1-y)\Lambda_\beta\bigr)
= -\log\!\left[y+(1-y)\frac{\Lambda_\beta}{\Lambda_\alpha}\right].
\label{eq:S-Lambda}
\end{equation}

For coherent-state excitations in the free boson CFT one may identify
\begin{equation}
\frac{\Lambda_\beta}{\Lambda_\alpha}=e^{-S(\alpha\|\beta)},
\qquad
S(\alpha\|\beta)=(\alpha-\beta)^2K(x),
\qquad
K(x)=1-\pi x\cot(\pi x).
\label{eq:Salpha}
\end{equation}
Substituting \eqref{eq:Salpha} into \eqref{eq:S-Lambda} gives the final answer
\begin{equation}
\boxed{
S\!\left(\rho_\alpha\,\Big\|\,y\rho_\alpha+(1-y)\rho_\beta\right)
=
-\log\!\left(
y+(1-y)\exp\!\big[-(\alpha-\beta)^2(1-\pi x\cot(\pi x))\big]
\right).
}
\label{eq:S-final}
\end{equation}

\section*{Consistency tests}
\label{sec:tests}

Let $s:=(\alpha-\beta)^2K(x)\ge 0$ and $a:=e^{-s}\in(0,1]$ so that
\begin{equation}
S(y)\equiv S\!\left(\rho_\alpha\,\Big\|\,y\rho_\alpha+(1-y)\rho_\beta\right)
=
-\log\!\lb y+(1-y)a\rb
=
-\log\!\lb a+y(1-a)\rb\,.
\end{equation}
(Notice the second equality is just algebra; the previous draft had a small typo here.)

We test the following limits:
\begin{itemize}
\item \textbf{Identical states:} $\alpha=\beta\Rightarrow s=0\Rightarrow a=1\Rightarrow S(y)=0$ as expected.
\item \textbf{No mixing:} $S(1)=0$ as expected and $S(0)=S(\alpha\|\beta)$ as expected. Therefore we can write
\begin{equation}
e^{-S(y)}=y\,e^{-S(1)}+(1-y)\,e^{-S(0)}\ .
\end{equation}
\item \textbf{Much heavier state:} $s\to\infty\Rightarrow a\to 0\Rightarrow S(y)\to -\log y$.
\end{itemize}

Next, we test the following properties:
\begin{itemize}
\item \textbf{Positivity:} since $y\le a+y(1-a)\le 1$, we have $0\le S(y)\le -\log y$ as required for relative entropy.
\item \textbf{Convexity in $y$:} taking derivatives with respect to $y$,
\begin{equation}
S'(y)= -\frac{1-a}{a+y(1-a)}\le 0,
\qquad
S''(y)=\frac{(1-a)^2}{(a+y(1-a))^2}\ge 0.
\end{equation}
Thus $S(y)$ is decreasing and convex in the mixing parameter $y$ (strictly convex if $s>0$),
consistent with the general theorem that $S(\rho\|\sigma)$ is convex in $\sigma$.

\item \textbf{Monotonicity and concavity in $s$:}
Treating $S$ as a function of $s$:
\begin{equation}
\frac{\partial S}{\partial s}=\frac{(1-y)e^{-s}}{y+(1-y)e^{-s}}\ge 0,
\qquad
\frac{\partial^2 S}{\partial s^2}
=-\frac{y(1-y)e^{-s}}{(y+(1-y)e^{-s})^2}\le 0.
\end{equation}
So $S$ increases with $s$ but is concave in $s$, consistent with saturation at $-\log y$.

\item \textbf{Monotonicity in region size:}
A direct derivative gives
\begin{equation}
K'(x)=\frac{\pi}{\sin^2(\pi x)}\left(\pi x-\frac12\sin(2\pi x)\right)>0
\qquad \text{for }x\in(0,1),
\end{equation}
since $\sin u<u$ for $u>0$. Thus $K(x)$ (and therefore $S$) is increasing in $x$ for $\alpha\neq\beta$, as expected from data processing.
\end{itemize}

Finally, we note that in the small interval limit $x\ll 1$ we have
\begin{equation}
K(x)=1-\pi x\cot(\pi x)=\frac{\pi^2}{3}x^2+\frac{\pi^4}{45}x^4+O(x^6).
\end{equation}
Hence
\begin{equation}
S(y)= (1-y)\,(\alpha-\beta)^2\frac{\pi^2}{3}x^2+O(x^4).
\end{equation}
The relative entropy vanishes as $x^2$ for small intervals, as expected.

\subsection*{Alternative analytic continuation}
\label{sec:small-n1}

An alternative analytic continuation follows from observing that since we are interested in small $n-1$, and the sum over $p=0,\cdots, n-1$, we effectively have $p=O(n-1)$ near $n=1$. Thus we linearize our functions in $p$ and discard higher order terms:
\begin{equation}
S_m(p)= -\frac{p}{2}+\frac{1-\cos(p\theta_m)}{4\sin^2(\theta_m/2)}
= -\frac{p}{2}+O(p^2).
\label{eq:Sm-smallp}
\end{equation}
Expanding the $g$ function in $n=1+\eps$ results in
\begin{equation}
\log g(p;x)=p\,L(n,x)+O(p^2)=p\,L(n,x)+O(\eps^2),
\end{equation}
uniformly over the relevant range $p\in\{0,\dots,n-1\}$. This reproduces the linear-in-$p$
structure that one would obtain by keeping only the ``smooth'' term $S_m(p)\approx -p/2$ and
effectively discarding the oscillatory $\cos(p\theta_m)$ piece at this order. While this results in a different analytic continuation than what we computed above, the $n\to 1$ limit gives the same result, and the same final expression \eqref{eq:S-final} for relative entropy.

\section*{Pole-matching analytic continuation of $g(p;x)$ and $h(n-p,p;x)$}
\NL{Check the discussion below }

In this appendix we reformulate the analytic continuation of the building blocks
$g(p;x)$ and $h(n-p,p;x)$ in the same spirit as the appendix of \texttt{arXiv:1508.03506}:
one writes a meromorphic candidate with the \emph{same poles (and the same pole orders)}
as the original expression at integer replica number, and then shows that the ratio is
bounded, hence constant by Liouville's theorem.


\


\section*{Information-theoretic interpretation}

% (placeholder)

\end{document}


\end{document}